% ===============================================================
%  Chapter 8 -- How the Cost Grows with Core Count       (Part II)
% ===============================================================
\chapter{How the Cost Grows with Core Count}
\label{ch:scaling}

\begin{figure}[htbp]
\centering
\adjustbox{max width=\linewidth}{%
\begin{forest}
  chapter map,
  [Chapter 8
    [Broadcast {\&} saturation [Lookups grow as N squared] [Bus demand grows with N]]
    [Point-to-point networks [Directory indirection]]
    [Coherence in AMAT [A fourth term]]
    [What this report targets [Not false sharing]]
  ]
\end{forest}}
\caption[]{Structure of Chapter~8. Each section shows one
way coherence cost grows with the number of cores; the last picks the one this
report studies.}
\label{fig:ch8-map}
\end{figure}

The models in this chapter are kept simple on purpose. They show how each cost
grows with the core count, and Part~VI gives measured values.

% ---------------------------------------------------------------
\section{Broadcast and Bus Saturation}
\label{sec:scale-broadcast}
\label{sec:scale-bus}
\sectionmap{[Broadcast cost [Each miss is snooped by N-1 caches]
  [Lookups grow as N squared] [Bus demand grows as N] [Saturation: requests queue]]}
\index{snooping protocol}\index{bus utilisation}%

On a snooping bus every other cache sees each miss and looks up its own tags to
decide whether to respond. With $N$ cores each making $r$ misses per unit time,
the snoop lookups per unit time are
\begin{equation}
  L(N) = N\,r\,(N-1),
  \label{eq:snoop}
\end{equation}
which grows with the square of the core count while useful work grows linearly
(\cref{fig:snoop-growth}). Each lookup costs energy and competes with the core's
own accesses for the tag array.

The bus also carries one message at a time (\cref{sec:coh-snoop}). If each core
keeps the bus busy for a fraction $d$ of the time, the utilisation with $N$ cores
is
\begin{equation}
  U(N) = N\,d
  \label{eq:util}
\end{equation}
until it reaches 100\,\%. Past that point the bus cannot carry the demand,
requests wait in the arbiter's queue and every miss takes longer
(\cref{fig:bus-util}). A core that needs the bus only 10\,\% of the time saturates
it at 10 cores, and every byte a protocol sends without need raises $d$ and brings
saturation closer.

\begin{figure}[htbp]
\centering
\begin{minipage}[t]{0.49\linewidth}\centering
\begin{tikzpicture}
  \begin{axis}[width=\linewidth, height=5cm, ybar, bar width=6pt,
      xlabel={Number of cores, $N$}, ylabel={Per unit time ($r=1$)},
      symbolic x coords={2,4,8,16}, xtick=data, ymin=0, ymax=290,
      legend style={at={(0.03,0.97)}, anchor=north west, font=\tiny},
      nodes near coords, nodes near coords style={font=\tiny},
      grid=major, grid style={coreGrey!15},
      label style={font=\footnotesize}, tick label style={font=\scriptsize}]
    \addplot[fill=tealC!60, draw=tealC] coordinates {(2,2) (4,4) (8,8) (16,16)};
    \addlegendentry{misses, $Nr$}
    \addplot[fill=amberC!70, draw=amberC!90!black] coordinates {(2,2) (4,12) (8,56) (16,240)};
    \addlegendentry{snoop lookups, $Nr(N-1)$}
  \end{axis}
\end{tikzpicture}
\caption[Snoop lookups grow quadratically]{\Cref{eq:snoop}: from 8 to 16 cores
the misses double and the lookups more than quadruple.}
\label{fig:snoop-growth}
\end{minipage}\hfill
\begin{minipage}[t]{0.49\linewidth}\centering
\begin{tikzpicture}
  \begin{axis}[width=\linewidth, height=5cm,
      xlabel={Number of cores, $N$}, ylabel={Bus demand $U(N)$ (\%)},
      xmin=1, xmax=16, ymin=0, ymax=200, xtick={2,4,8,12,16},
      legend style={at={(0.03,0.97)}, anchor=north west, font=\tiny},
      grid=major, grid style={coreGrey!15},
      label style={font=\footnotesize}, tick label style={font=\scriptsize}]
    \addplot[fill=alertRed!10, draw=none, forget plot] coordinates {(1,100) (16,100) (16,200) (1,200)} \closedcycle;
    \addplot[very thick, leafC, domain=1:16] {5*x};   \addlegendentry{$d=5\,\%$}
    \addplot[very thick, plumC, domain=1:16] {10*x};  \addlegendentry{$d=10\,\%$}
    \addplot[very thick, brickC, domain=1:16] {20*x};  \addlegendentry{$d=20\,\%$}
    \addplot[very thick, dashed, coreGrey, domain=1:16, forget plot] {100};
    \node[font=\tiny, text=alertRed] at (axis cs:11,185) {demand exceeds capacity};
  \end{axis}
\end{tikzpicture}
\caption[Bus saturation]{\Cref{eq:util}. Above the dashed line requests queue;
lowering $d$ moves saturation to the right.}
\label{fig:bus-util}
\end{minipage}
\end{figure}

% ---------------------------------------------------------------
\section{Point-to-Point Networks and Directories}
\label{sec:scale-p2p}
\sectionmap{[Beyond the bus [Mesh carries many messages at once]
  [Directory avoids broadcast] [Price: an extra hop]]}
\index{MESI}\index{snooping protocol}\index{directory protocol}\index{network on chip}%

Larger chips replace the bus with a point-to-point network such as a mesh
(\cref{tab:topologies}) and use a directory in place of broadcast
(\cref{sec:coh-dir}). The quadratic snoop cost disappears and many messages can
travel at once, but coherence still costs something. A read miss to a block
modified in another cache now needs three network traversals: request to the
directory, forward to the holder, and reply to the requester. Each crosses several
links of a network whose diameter grows with
$\sqrt{N}$~\cite[\S12.6.3]{sarangi2025basic}. Under MESI the holder also writes
the block back, so the block still moves twice (\cref{ch:scaling-approaches}
draws this transaction).

% ---------------------------------------------------------------
\section{Coherence in the AMAT Equation}
\label{sec:scale-amat}
\sectionmap{[AMAT with coherence [Add a coherence-miss term]
  [Rate grows with sharing and cores] [Cost includes any write back]]}
\index{AMAT}\index{coherence miss}%

\Cref{eq:amat} covers the three kinds of miss of a single core. Adding coherence
misses (\cref{sec:orig-4c}) gives a fourth term,
\begin{equation}
  \mathrm{AMAT} = T_{L1} + m_{L1}\bigl(T_{L2} + m_{L2}\,T_{\mathrm{mem}}\bigr)
                + m_{\mathrm{coh}}\,T_{\mathrm{coh}},
  \label{eq:amat-coh}
\end{equation}
where $m_{\mathrm{coh}}$ is the fraction of accesses that are coherence misses and
$T_{\mathrm{coh}}$ is the time to get the block from another cache, including any
writeback on the way. \Cref{fig:amat-bars} applies \cref{eq:amat-coh} to the
3-cycle example of \cref{sec:cache-hierarchy} with an assumed
$T_{\mathrm{coh}}=60$ cycles and a coherence-miss rate that doubles each time the
core count doubles. Under these assumptions the coherence term is larger than the
whole single-core hierarchy by 16 cores.

\begin{figure}[htbp]
\centering
\begin{tikzpicture}
  \begin{axis}[width=9.5cm, height=4.6cm, ybar stacked, bar width=20pt,
      xlabel={Number of cores (assumed $m_{\mathrm{coh}}$)}, ylabel={AMAT (cycles)},
      symbolic x coords={a,b,c,d}, xtick=data, xticklabels={{2 (1\%)},{4 (2\%)},{8 (4\%)},{16 (8\%)}},
      ymin=0, ymax=9,
      legend style={at={(0.03,0.97)}, anchor=north west, font=\scriptsize},
      grid=major, grid style={coreGrey!15},
      label style={font=\small}, tick label style={font=\scriptsize}]
    \addplot[fill=skyC!55, draw=skyC] coordinates
      {(a,3) (b,3) (c,3) (d,3)};
    \addlegendentry{single-core terms}
    \addplot[fill=roseC!55, draw=roseC] coordinates
      {(a,0.6) (b,1.2) (c,2.4) (d,4.8)};
    \addlegendentry{coherence term $m_{\mathrm{coh}}T_{\mathrm{coh}}$}
  \end{axis}
\end{tikzpicture}
\caption[AMAT with a coherence term]{\Cref{eq:amat-coh} with illustrative values.
A smaller $T_{\mathrm{coh}}$, for example one without a writeback on every
coherence transfer, shrinks the pink bars.}
\label{fig:amat-bars}
\end{figure}

% ---------------------------------------------------------------
\section{What This Report Targets}
\label{sec:scale-target}
\label{sec:scale-false}
\sectionmap{[Target [Many sources of coherence cost]
  [Most are inherent to sharing] [{False sharing: not targeted}] [One is avoidable in principle]]}
\index{MESI}\index{MSI}\index{false sharing}%

Coherence works on blocks. If two cores write different variables that happen to
share a block, each write invalidates the other core's copy and the block moves
back and forth even though no data is shared. This false
sharing\index{false sharing} causes real traffic, but it comes from writes on both
sides, and every MSI-family protocol must transfer ownership on each of them, so
we leave it out.

\Cref{fig:overhead-map} collects the costs from Chapters~7 and~8. Most come with
sharing itself: a block has to move when another core needs it, and other copies
have to be invalidated when a core writes. The last one is different. When a core
reads a block that another core holds modified, MESI sends the block to the
reader, which is required, and also writes it back to memory, which no rule of
coherence requires. Chapter~9 looks at that writeback.

\begin{figure}[htbp]
\centering
\adjustbox{max width=0.85\linewidth}{%
\begin{forest}
  for tree={
    grow'=0, parent anchor=east, child anchor=west, anchor=west,
    align=center, font=\sffamily\small, draw, thick, rounded corners=3pt,
    inner sep=4pt, l sep+=14pt, s sep+=2pt,
    edge={thick, coreGrey!60, -{Stealth[length=5pt]}},
    fill=coreGrey!8, draw=coreGrey!60, text=coreGrey
  },
  [Coherence cost, fill=mesiBlue, draw=mesiBlue, text=white, font=\sffamily\bfseries\small
    [Snoop lookups on every miss]
    [Bus or network contention]
    [Invalidations on writes]
    [Cache-to-cache data transfers]
    [False sharing]
    [Write backs when a modified\\block is read by another core,
      fill=moesiOrange, draw=moesiOrange, text=white, font=\sffamily\bfseries\small
      [{Targeted by\\this report}, fill=moesiOrange!12, draw=moesiOrange, text=black]]
  ]
\end{forest}}
\caption[Sources of coherence cost]{Sources of coherence cost found in
Chapters~7 and~8. The highlighted one is studied in this report.}
\label{fig:overhead-map}
\end{figure}
